\begin{textsample}{POS Dim 1 – al – Score 40.00 – al/al\_ef\_55.txt}  \label{ex:f1_pos_152}
2.
CONCEPÇÕES HISTÓRICAS E EVOLUÇÃO DO COMPONENTE No Brasil, o ensino de \textbf{uma} língua estrangeira, em especial da língua inglesa, se pauta \textbf{historicamente} na trilogia : língua, cultura e identidade, \textbf{uma} vez \textbf{que} esses aspectos são fundamentais na práxis pedagógica dos professores da língua inglesa.
Esses conceitos estão interligados nessa visão histórica do ensino/aprendizagem dessa língua no cenário da educação básica brasileira.
A língua é o principal foco dos estudos das ciências linguísticas ; entretanto, esse termo dispõe de variados conceitos e é fundamental apresentá -los, ainda \textbf{que} sucintamente, \textbf{visto} \textbf{que}, numa perspectiva histórica, ela assume diferenciados papéis nas sociedades.
No início do \textbf{século} XX, Ferdinand de Saussure criou a Linguística Estruturalista, cuja \textbf{abordagem} tem como aspecto mais relevante e importante a definição da língua como objeto de estudo.
Saussure faz \textbf{uma} distinção de linguagem na qual afirma \textbf{que} a linguística é constituída por todas as manifestações humanas e \textbf{que} tem duas \textbf{divisões} : Língua — considerada essencial — e Fala, \textbf{vista} como secundária.
Assim : > “ O estudo da linguagem comporta, portanto, duas partes : \textbf{uma}, essencial, tem por objeto a língua, \textbf{que} é social em sua essência e independente do indivíduo ; esse estudo é unicamente psíquico ; outra, secundária, tem por objeto a parte individual da linguagem, vale \textbf{dizer}, a fala, inclusive a fonação e é psicofísica ( SAUSSURE, 2000, p. 27 ). ” O referido \textbf{autor} entende \textbf{que} a língua é \textbf{um} sistema homogêneo, imutável ao longo do tempo e cuja existência funda -se na necessidade de comunicação, entre outros fatores.
A visão Saussuriana de língua não considera a variação linguística, ou seja, a língua “ é \textbf{um} fato social porque pertence a todos os membros de \textbf{uma} comunidade, é exterior ao indivíduo, esse não pode nem criá -la nem modificá -la ” ( CARDOSO, 1999, p. 15 ). \textbf{Contemporâneo} de Saussure, Mikhail Bakhtin mostra outra visão de língua, a qual é \textbf{amplamente} estudada como suporte \textbf{teórico} no ensino de línguas.
Para este, a língua \textbf{vive} e evolui \textbf{historicamente} na comunicação verbal concreta.
É, ainda, \textbf{um} fenômeno essencialmente mutável, social, \textbf{que} ocorre na interação entre os indivíduos e está em constante evolução, recriação e possibilidade de modificação.
Para ambos, a concepção de língua é \textbf{vista} como fator social.
No entanto, Bakhtin prioriza a questão dos dados reais da linguística, ou seja, estuda a natureza social da linguagem, a fala do indivíduo e seu caráter intencional.
O \textbf{teórico} contrapõe -se ao objetivismo abstrato de Saussure em relação à língua, mostra \textbf{que} as formas dela são \textbf{um} produto de relações sociais pelos interlocutores e não definidas exclusivamente por \textbf{um} sistema abstrato de formas linguísticas — concepção defendida por Saussure no Estruturalismo.
Levando -se em consideração a visão de Bakhtin ( 1997 ), no tocante à concepção da língua como \textbf{um} processo \textbf{dialógico}, social e de interação verbal, é \textbf{preciso} \textbf{atentar} -se ao ensino de língua estrangeira a partir dessa \textbf{abordagem} \textbf{dialógica} bakhtiniana — entre outras —, \textbf{uma} vez \textbf{que}, nela, o ensino da LE e, em especial, a Língua Inglesa possibilita \textbf{uma} aprendizagem pautada em situações reais e cotidianas em \textbf{que} o \textbf{aprendiz} a percebe como importante e relevante para sua formação cidadã.
Entretanto, apesar dessa visão já vir sendo difundida por todo esse tempo, muitos professores ainda tendem, por razões e fatores diversos, a dar \textbf{um} tratamento ao ensino da língua inglesa pautado numa visão estrutural, com ênfase nos aspectos gramaticais em detrimento aos demais aspectos linguísticos fora do contexto social, cultural e interacional, opondo -se, por \textbf{conseguinte}, à visão da língua defendida na BNCC, como \textbf{uma} língua de interação social, multicultural e franca.
Fruto de \textbf{uma} construção histórica, na educação básica brasileira, há vários documentos \textbf{que} subsidiam o caminhar dos educadores em seus componentes curriculares rumo a \textbf{uma} proposta de ensino pautada no dinamismo da sociedade \textbf{atual} \textbf{que} se depara com \textbf{um} mundo plural e multidimensional.
Dentre esses, vale destacar a fundamental importância das Orientações Curriculares para o Ensino Médio ( Brasil, 2006 ), cuja função básica é orientá -los na atuação em sala de aula, além de contribuir para o diálogo envolvendo a tríade — o professor, sua prática docente e a escola —, sendo \textbf{um} instrumento de apoio \textbf{teórico} ao professor a ser utilizado em favor do aprendizado.
Nesse material, encontra -se manifestada, também, \textbf{uma} posição em relação à concepção de língua, ao afirmar \textbf{que} : > [... ] é na interação em diferentes instituições sociais ( a família, o grupo de amigos, as comunidades de bairro, as igrejas, a escola, o trabalho, as associações, etc. ) \textbf{que} o \textbf{sujeito} apreende as formas de funcionamento da língua e os modos de manifestação da linguagem ; ao fazê -lo, vai construindo seus conhecimentos relativos aos usos da língua e da linguagem em diferentes situações.
Também nessas \textbf{instâncias} sociais o \textbf{sujeito} constrói \textbf{um} conjunto de representações sobre o \textbf{que} são os sistemas semióticos, o \textbf{que} são as variações de uso da língua e da linguagem, bem como qual seu valor social ( BRASIL, 2006, p. 24 ).
A visão proposta pelas OCEM, \textbf{que}, apesar de referir -se ao ensino médio, é de fundamental relevância \textbf{aqui}, \textbf{visto} \textbf{que} norteará a concepção de língua para o ensino/aprendizagem do estudante egresso do ensino fundamental, reforça a visão Bakhtiniana de \textbf{uma} língua construída pela e na interação social e histórica de seus falantes/usuários.
É nesse viés de língua como interação entre \textbf{sujeitos} \textbf{que} se faz necessário destacar a importância da Linguística Aplicada e sua contribuição nessa concepção \textbf{enquanto} função social na aprendizagem de língua estrangeira.
A referida ciência surgiu na segunda metade do \textbf{século} XX como a aplicação da linguística quando ela ampliava suas \textbf{abordagens} históricas.
No entanto, ao longo do tempo, passou a dialogar com a própria Linguística e outras \textbf{disciplinas}, tornando -se área de natureza heterogênea, comprometida com questões político-econômicas e também voltadas para o uso da linguagem na sala de aula.
Em contraste com o modelo estruturalista da língua, a LA tem se apresentado numa perspectiva de língua real, pautada na prática social, \textbf{que} muda com a evolução linguística, como afirma Signorini ( 1998, p. 101 ) : > A LA tem buscado cada vez mais a referência de \textbf{uma} língua real, ou seja, \textbf{uma} língua falada por falantes reais em suas práticas reais e específicas, numa tentativa justamente de seguir essas redes, de não arrancar o objeto da tessitura de suas raízes.
Pela \textbf{exposição} das concepções de língua acima, fica claro \textbf{que} tais visões são bem distintas entre si e \textbf{que} são relevantes e coerentes em seus momentos de estudo e contexto, pois, quando analisadas e entendidas em seus momentos, “ percebe -se \textbf{que} a evolução da língua e suas transformações históricas nas sociedades complexas, dinâmicas e em constante evolução requerem mudanças \textbf{que}, por sua vez, incidirão na representação da linguagem e cabe ao professor estar atento a essa nova representatividade ” ( TEIXEIRA ; RIBEIRO, 2013, p. 119 ).
Quando se pensa numa língua falada por \textbf{um} povo, essa língua envolve outros aspectos como cultura, etnia, \textbf{raças} para além de \textbf{fronteiras}.
É na língua(gem) \textbf{que} o \textbf{homem} constrói sua existência dentro de certas condições socioeconômicas \textbf{que} \textbf{desencadeiam} sua produção cultural, sua representatividade e suas necessidades humanas de se \textbf{agregar} socialmente, de se mostrar como indivíduo.
Quando se estuda \textbf{uma} língua estrangeira, \textbf{aqui} com ênfase à língua inglesa, é de \textbf{suma} importância discutir o(s) sujeito(s) \textbf{aprendiz} e sua(s) identidade(s) como sujeito(s) cultural em suas múltiplas concepções, \textbf{que} (re)conhece o outro e convive real e ou virtualmente com outras culturas constituídas na heterogeneidade.
O contexto escolar é \textbf{um} vasto e fértil campo para construção e ampliação dessas relações. \textbf{Uma} vez \textbf{que} é nesse espaço e, em especial na sala de aula, \textbf{que} o estudante amplia seu contato social e vai construindo no decurso das etapas escolares : amizades, conhecimentos, aprendizagens e costumes ; e, dessa \textbf{vertente}, a língua inglesa tem tido como \textbf{um} dos objetivos oportunizar ao \textbf{aprendiz} conhecer e apreender essa nova língua para fruição, informação, compreensão de outros modos de vida, desenvolvendo com criticidade sua visão sobre as desigualdades sociais ( de gênero, \textbf{raça}, classe, \textbf{sexualidade}, dentre outros ).
Com a evolução da educação através dos tempos e da língua inglesa no cenário nacional, fica nítido o entendimento de \textbf{que} a língua inglesa deve levar em conta, por parte do professor e da escola, o seu sujeito/estudante em seu caráter individual, social, cultural no espaço do \textbf{ensino-aprendizagem} e propiciar a esse \textbf{sujeito} ser agente de seu próprio conhecimento, levando -o a \textbf{uma} aprendizagem efetiva, interativa, reflexiva a respeito de si e o do outro, reconhecendo -se como \textbf{um} \textbf{sujeito} social inacabado e em constante transformação e preparação para inserir -se e sentir -se cidadão num mundo \textbf{globalizado}.
Em Alagoas, quando há referência ao ensino e aprendizagem de língua inglesa nas escolas do ensino fundamental e médio, muitas vezes há \textbf{um} olhar de desconfiança.
Dentre as várias justificativas, \textbf{uma} é de \textbf{que} é muito difícil e complicado para o estudante conhecer a língua inglesa suficientemente bem com a carga horária oferecida, sendo necessários tempo e \textbf{bastante} acúmulo de conhecimento para \textbf{que} se possa, ao final, falar ou escrever com fluência nessa língua.
Outra justificativa é \textbf{que} as escolas não dispõem de estrutura suficiente e adequada à aprendizagem e \textbf{que} os estudantes não se interessam pela língua inglesa porque ela está muito distante da realidade deles.
Não obstante esses fatos, ainda há o despreparo de alguns profissionais, o \textbf{que} os tem levado a \textbf{uma} prática pedagógica centrada na língua em seu aspecto meramente estruturalista e obsoleto, \textbf{que} em nada atrai o estudante e \textbf{que} contraria as novas propostas da BNCC para esse componente.
Chega -se, assim, à conclusão de \textbf{que} a escola pública não é lugar para se aprender a língua inglesa.
Chama -se a isso de mito do fracasso do ensino de inglês na escola pública.
Ao tratarem desse assunto, as Orientações Curriculares ( OCNEM ) para o Ensino Médio ( BRASIL, 2006 ) informam \textbf{que} esse mito se constitui no motor da proliferação de cursos livres no país.
Isso cria \textbf{um} fator de \textbf{exclusão} violento porque aqueles \textbf{que} podem pagar e, portanto, estudar inglês, poderão \textbf{conseguir} melhores empregos, acessar informações em locais distantes e interagir com o outro.
No Estado de Alagoas, esse mito também é recorrente.
Em entrevistas com professores de inglês das escolas públicas estadual, por ocasião da construção do referencial curricular da rede, a maioria reconhece a importância da língua inglesa nos currículos escolares, mas é unânime em \textbf{dizer} \textbf{que} o estudante alagoano da escola pública vem de \textbf{um} contexto muito pobre e, por isso, não tem condições de aprender inglês ou outra língua ofertada, como o espanhol, por exemplo.
Paradoxalmente, os estudantes, por sua vez, utilizam o computador, telefone celular com acesso à internet ; assistem televisão ; vão às festas ; participam de encontros sociais e religiosos ; interagem com seus amigos, \textbf{conhecidos} \textbf{face} a \textbf{face} e por meio virtual, etc.
Esses contextos em \textbf{que} os estudantes estão inseridos são muito ricos culturalmente e oferecem oportunidades de interação contínua em língua inglesa também, o \textbf{que} contraria a ideia de \textbf{que} o estudante da escola pública estaria fadado ao “ mito do fracasso ”.

% matched lemmas: abordagem, agregar, amplamente, aprendiz, aqui, atentar, atual, autor, bastante, conhecido, conseguinte, conseguir, contemporâneo, desencadear, dialógico, disciplina, divisão, dizer, enquanto, ensino-aprendizagem, exclusão, exposição, face, fronteira, globalizado, historicamente, homem, instância, preciso, que, raça, sexualidade, sujeito, sumo, século, teórico, um, ver, vertente, viver
\end{textsample}
